\begin{definition}[Semiconvexity]
Let $E$ be a finite-dimensional real inner-product space, let $\lambda\in\R$, let
$\Omega\subseteq E$ and let $\phi:E\to\R$. We say that $\phi$ is
\emph{$\lambda$-semiconvex on $\Omega$} if the function
\[
  z\longmapsto\phi(z)+\frac{\lambda|z|^2}{2}
\]
is convex on $\Omega$. Here, as throughout, ``$g$ is convex on $\Omega$'' means
that $\Omega$ is a convex subset of $E$ and that
\[
  g(t_1z_1+t_2z_2)\le t_1g(z_1)+t_2g(z_2)
\]
for all $z_1,z_2\in\Omega$ and all $t_1,t_2\ge0$ with $t_1+t_2=1$.

Thus $0$-semiconvexity is convexity, and for $\lambda>0$ a
$\lambda$-semiconvex function is one that is convex after the addition of the
quadratic $\lambda|z|^2/2$; informally its second derivative, where it exists,
is bounded below by $-\lambda$ times the identity.
\end{definition}
